\documentclass[12pt]{article}
\usepackage[utf8]{inputenc}
\usepackage[T1]{fontenc}
\usepackage{hyperref}
\title{Pontic Language Resources for K-12 Curriculum Development}
\author{Compiled by Claude Sonnet 5.5 (claude-sonnet-5-5) for Andr\'as Kornai}
\date{}
\begin{document}
\maketitle

\section{Introduction}
Pontic (Pontic Greek, Romeika, Romeyka; ISO 639-3 pnt) is a Greek variety of the Black Sea region, spoken by communities in Greece and Turkey. This report lists only resources that appeared in web search results. The search found no Pontic school curriculum and no Pontic--English learner's course; the material below is mostly lexical and descriptive.

\section{Primary Educational Resources}
No primers, textbooks or school programmes were found in the search results. PontosWorld is described as offering lessons and introductory articles within its dialect section, but their content and suitability for schools were not verified.
\begin{itemize}
\item PontosWorld, dialect overview (dictionary, lessons and articles; not verified in detail): \url{https://pontosworld.com/index.php/dialect/overview}
\end{itemize}

\section{Linguistic and Reference Materials}
\begin{itemize}
\item Vahit Tursun, \emph{Romeika--T\"urk\c{c}e S\"ozl\"uk} (Romeika--Turkish dictionary, Heyamola, 2019), about 14,400 entries, with an introduction by Peter Mackridge; catalogue entry: \url{https://mail.pontosworld.com/index.php/books/pontic/847-romeika-turkce-sozluk}
\item University of Bamberg repository record of a thesis described as the first thorough grammatical description of Romeyka, the Muslim variety of Pontic Greek in Trabzon province: \url{https://fis.uni-bamberg.de/handle/uniba/95115}
\item Languagehat post on a Pontic Greek dictionary (discussion, not a primary resource): \url{https://languagehat.com/pontic-greek-dictionary/}
\end{itemize}

\section{Religious and Community Resources}
No scripture translations or community organisations were identified in the search results beyond the PontosWorld site listed above.

\section{Digital and Audio Resources}
The Endangered Languages Project lists alternate names and dialect divisions for the language, but its speaker and support fields were empty:
\begin{itemize}
\item \url{https://endangeredlanguages.com/node/96251}
\end{itemize}
No audio or video collections were found.

\section{Curriculum Development Context}
Documentation exists, but the search found no school materials. A curriculum would have to be built from the dictionary and grammar above, together with speaker work. The Bamberg description reports ongoing language shift to Turkish among Romeyka speakers, so documentation work is time-sensitive. Further searching (Greek-language sources, Pontic cultural associations) is needed.
\end{document}